% Notes of the Montreal seminar, 1970 (cote 129, typed part): archive cote 129, site page 61 (PDF page 62).
% typescript. Draft compared: none.
% AI-transcribed by Claude Opus 5.5 (claude-opus-5-5, adaptive thinking, effort high), Message Batches API, batch mtl129_full (msgbatch_01AGnDKHwy6ATCM7ehUa6263), 2026-11-13; not yet checked by eye.
\scanpage{129/61}
\mspage{III-6}

3.2. \uline{Définition}. Soit $n \geq 1$ un entier. On appelle \uline{p-groupe de}
\uline{Barsotti-Tate tronqué d'échelon $n$ sur $S$} (ou groupe Barsotti-Tate
tronqué d'échelon $n$) un schéma en groupes (commutatifs) $G$ fini et
localement libre sur $S$, annulé par $p^{n}$ et plat sur $\Lambda_{n}$. Si $n = 1$,
on suppose de plus que l'on a (avec les notations de 3.1):
\[
(*) \qquad \operatorname{Ker} \mathbb{F}_{G_{o}/S_{o}} = \operatorname{Im} \mathbb{V}_{G_{o}}
\]
(c'est une condition qui se vérifie fibre à fibre).

3.3. Si $G(n)$ est un groupe de BT tronqué d'échelon $n$, alors
$G(i)$, $1 \leq i \leq n$, est un groupe de BT tronqué d'échelon $i$.

Pour $i > 1$, cela résulte de 2.4; pour $i = 1$, nous allons
montrer que la condition (*) est aussi vérifiée.

Au dessus de $S_{o} \subset S$, nous avons déjà remarqué que l'image de
$\mathbb{V}_{G(1)_{o}} : G(1)_{o}^{(p)} \rightarrow G(1)_{o}$ était contenue dans le noyau $K$ de
$\mathbb{F}_{G(1)_{o}/S_{o}} : G(1)_{o} \rightarrow G(1)_{o}^{(p)}$ , de sorte que l'on a un diagramme commutatif
\diagram{
\begin{tikzcd}
G(1)_{o}^{(p)} \arrow[r, "\mathbb{V}"] \arrow[dd] & K \arrow[d, hook] \\
 & G(1)_{o} \arrow[d, hook] \\
G(2)_{o}^{(p)} \arrow[r, "\mathbb{V}_{G(2)_{o}}"] & G(2)_{o}
\end{tikzcd}
}
dont on vérifie immédiatement qu'il est cartésien car $G(2)_{o}^{(p)}(1) \simeq G(1)_{o}^{(p)}$.
Par ailleurs l'image de $\mathbb{V}_{G(2)_{o}}$ contient $G(1)_{o}$ , car
$\mathbb{V}_{G(2)_{o}} \circ \mathbb{F}_{G(2)_{o}} = p$ a pour image $G(1)_{o}$ , et l'on en déduit que
$\mathbb{V} : G(1)_{o}^{(p)} \rightarrow K$ est surjectif.

\mspage{61}
